% 第五章示例：总结主要成果并提出后续研究方向。
\chapter{总结与展望}

\section{总结}

最后一章应着重总结毕业设计（论文）的创新点、新见解和主要研究
成果。总结是对论文主要工作、结果和论点的提炼与概括，应准确、
简明、完整且有条理，使读者能够全面了解论文的研究意义、目的和
工作内容。

总结应严格区分作者本人取得的成果与指导教师及其他研究者的成果。
评价研究工作时必须实事求是；除非有充分证据，不宜使用“首次”、
“领先”或“填补空白”等绝对化表述。

\section{展望}

展望或建议应建立在已完成工作和现有结论的基础上，对相关技术或研究
领域未来的发展方向作出合理预测，并评价研究成果的应用前景和社会
影响，为后续研究提供可借鉴的思路。
